\section{Introduction}
\label{sec:introduction}

Submarine cables form the physical backbone of the global Internet and carry
nearly all intercontinental traffic~\cite{telegeography_submarine_map}.
For users accessing services across oceans, these cables provide the physical
connections underlying end-to-end network paths. Traceroute exposes the
network-layer structure of these paths through visible routers and
Autonomous System (AS) transitions between sources and service destinations.

Failures in this physical infrastructure can isolate regions, degrade
remote-service access, and redirect traffic toward longer or congested
alternatives. Paths traversing different routers or ASes may nevertheless
share a cable or coastal corridor and be affected by the same physical event.
Such shared dependencies are not directly visible in traceroute: a service
may present many distinct network transitions while its feasible physical
support remains concentrated in a few coastal directions. Assessing
service-path diversity therefore requires comparing how observations are
distributed across network transitions and feasible physical corridors.

Cross-layer mapping provides a starting point for this comparison.
Internet Atlas, InterTubes, and iGDB connect observed Internet links to
long-haul infrastructure; Nautilus associates IP links with candidate
submarine cables; Calypso maps routes using cable layouts, network
relationships, and inland connectivity; and Xaminer uses cross-layer maps
to analyze infrastructure risk~\cite{durairajan2013atlas,durairajan2015intertubes,
anderson2022igdb,ramanathan2023nautilus,wang2025calypso,
ramanathan2024xaminer}. These studies make physical dependencies accessible
to measurement and analysis. Building on this capability, we ask whether
the diversity visible in a population of service paths reflects how its
feasible physical support is distributed.

Answering this question requires more than counting distinct network
transitions or feasible corridors. Different transitions may repeatedly
support the same corridors, while a recurring transition may admit several
physical candidates. Counts alone do not describe how observations are
distributed among those categories. Comparing the two layers therefore
requires retaining each segment's network identity and physical candidate
set, and aggregating both over the same service-path population.

We present \emph{CLASP}, a Cross-Layer Audit of Service Paths, to construct
this paired comparison. CLASP extracts atomic segments from traceroutes,
records their network-transition labels, and identifies feasible
landing-region corridors using geographic, cable-lifecycle, and propagation
constraints. Each candidate-bearing segment contributes the same total
observation mass to both layers; candidate allocation distributes its
physical-layer mass among feasible corridors. Aggregating by country and
measurement then produces paired distributions while preserving the
source-country and measured-target context. This construction compares
the organization of a fixed set of observations rather than distributions
drawn from different path populations.

The paired distributions support two complementary analyses. Within a
country--measurement unit, we measure changes in concentration and breadth
between network transitions and feasible corridors. Across units, we measure
whether network concentration preserves the ordering of corridor
concentration, using Spearman rank correlation, termed \emph{Layer Agreement}.
These analyses address \textbf{RQ1}, how diversity differs within units, and
\textbf{RQ2}, how concentration-rank association varies by geography and
candidate allocation. \textbf{RQ3} examines exposure, absolute corridor
breadth, and measurement-family composition as context for these results.
Exposure is reported separately because the frequency of encountering a
feasible corridor and the distribution of support among corridors describe
different properties.

We apply CLASP to 490,911 valid IPv4 traceroutes from 18 public RIPE Atlas
measurements covering all 13 DNS Roots, Wikipedia, Reddit, Netflix Assets,
and two dynamic multi-target topology references. The aligned one-hour
snapshot yields 370 auditable country--measurement units. The service
measurements retain their selected targets, while the multi-target
measurements provide a broader-target reference. Equal-share,
projection-score, and Top-1 allocation compare how the results depend on
distributing observation mass among feasible candidates.

The audit identifies service-facing populations whose broad network
support accompanies concentrated corridor support, as well as narrowing
within populations classified as broad at both layers. Across geographic
groups, concentration rankings track more closely among island and
archipelagic units than among coastal-mainland units under all three
allocation rules. The two findings describe different aspects of the
cross-layer relationship: aligned rankings can coexist with within-unit
narrowing. Together, they motivate examining physical-corridor distributions
alongside network-layer summaries when characterizing service-path diversity.

This paper makes three contributions. First, CLASP constructs paired
network-transition and feasible-corridor distributions over identical
traceroute segments, separating trace-level exposure from conditional
corridor concentration. Second, the audit measures within-unit diversity
differences and geographic variation in concentration-rank association,
with three candidate-allocation rules providing complementary views of
the feasible support. Third, it records measurement IDs, configuration
values, mapping states, and pipeline counts so that the reported
distributions can be reconstructed and inspected.
